\section{Introduction}

Evidence-based fact-checking asks a system to justify a verdict with identifiable sources rather than infer veracity only from the wording or appearance of a claim. The requirement becomes harder when both claims and evidence are multimodal: a decisive statement may occur in a paragraph, while an accompanying image supplies complementary or contradictory information. MOCHEG formalizes this setting as joint text--image evidence retrieval, three-way claim verification, and explanation generation \cite{yao2023mocheg}. Later systems add dense retrieval, image--text verbalization, language-model verification, or graph-based evidence reasoning \cite{luu2027mrav,zhang2026mever}. These advances make evidence selection itself a measurable part of the task.

A correct verdict and a useful evidence set are related but different outputs. A model can predict the right label while citing irrelevant items, omitting a necessary image, or returning many redundant documents. Generated explanations can also be fluent without accurately attributing their statements to the selected evidence \cite{xing2025attribution,niu2024ragtruth}. Evaluation must therefore preserve the identity of selected items and measure their overlap with released gold evidence instead of treating a verdict or plausible explanation as sufficient.

Most multimodal retrieval systems expose a globally fixed top-$k$ set. A single budget is simple and reproducible, but it assumes that every claim needs the same number of items. Choosing a larger $k$ can recover more gold evidence while reducing precision; choosing a smaller $k$ has the opposite effect. Predicting only the gold count does not solve this trade-off when the ranking is imperfect. The preceding SciVer study in this project isolated citation-set size from verification and obtained a negative result: a learned count head improved count accuracy but did not beat the strongest fixed budget in evidence-set F1. That consumed benchmark is retained only as motivation and is not used to select or confirm the new method.

We study whether the selection rule can target set utility directly. Calibrated Expected-F1 Evidence Selection (\method) estimates the probability that each candidate belongs to the gold evidence set, orders candidates by that probability, and evaluates every admissible prefix. An exact Poisson-binomial dynamic program computes the expected set-F1 for each prefix under an explicitly stated independent-Bernoulli approximation. The decoder returns the smallest prefix attaining the maximum expectation. Fixed-budget, threshold, count-prediction, and \method{} outputs use the same candidate packet, frozen features, and scorer family; the planned end-to-end comparison also uses the same verifier checkpoint. This design makes the selector, rather than model capacity or retrieval access, the controlled variable.

The study uses only public MOCHEG data. Train and development transports are pinned to immutable revisions and reconciled with official counts. The official test transport is recorded but its annotations remain unavailable to development code. A 400-claim feasibility packet has passed metadata, sentence-recovery, media-readability, gold-availability, and deterministic-construction checks. The implemented software includes the exact decoder, deterministic temperature calibration, pinned BGE-M3 and SigLIP2 feature extraction, fixed-budget controls, three development seeds, and paired bootstrap analysis. The GPU semantic pilot and matched-verifier experiment remain pending, so this manuscript is a pre-results study draft rather than a submission-ready empirical paper.

The intended contributions, conditional on the prespecified evaluation, are:

\begin{itemize}
    \item an adaptive multimodal evidence selector that optimizes exact expected set-F1 rather than citation count or a global retrieval depth;
    \item a matched comparison protocol that shares candidates, representations, scorer inputs, tuning access, and verifier checkpoints across adaptive and fixed selectors;
    \item a test-sealed MOCHEG evaluation with calibration, evidence-count, modality, duplicate, and robustness analyses; and
    \item an auditable development gate that prevents the official test from being consumed when the proposed selector does not establish a meaningful development advantage.
\end{itemize}

The paper does not claim that these contributions improve performance until the real semantic run, verifier non-inferiority test, frozen formal evaluation, and independent recomputation audit are complete.
